% =====================================================================
%  TEMPLATE TESIS — BERKAS INDUK
%  Kompilasi: pdflatex -interaction=nonstopmode main.tex  (jalankan 2x)
%  Struktur:
%    main.tex                    ← berkas ini (preamble + \include)
%    frontmatter/abstrak.tex     ← abstrak + kata pengantar
%    chapters/bab1–bab7.tex      ← bab-bab tesis
%    backmatter/referensi.tex    ← daftar pustaka
%    assets/                     ← gambar (.png/.pdf)
% =====================================================================

\documentclass[12pt,a4paper]{report}

% ── Bahasa & font ────────────────────────────────────────────────────
\usepackage[T1]{fontenc}
\usepackage[utf8]{inputenc}
\usepackage[indonesian]{babel}   % otomatis: Bab / Gambar / Tabel / Daftar Isi
\usepackage{mathptmx}            % Times New Roman
\usepackage{setspace}
\onehalfspacing                  % spasi 1,5

% ── Tata letak halaman ───────────────────────────────────────────────
\usepackage[a4paper,left=4cm,right=3cm,top=3cm,bottom=3cm]{geometry}
\usepackage{parskip}
\setlength{\parindent}{1.27cm}
\setlength{\parskip}{0pt plus 1pt}

% ── Tabel & gambar ───────────────────────────────────────────────────
\usepackage{graphicx}
\graphicspath{{assets/}}
\usepackage{booktabs}
\usepackage{array}
\usepackage{multirow}
\usepackage{longtable}
\usepackage{float}
\usepackage{caption}
\captionsetup{font=small, labelfont=bf, justification=centering}

% ── Matematika ───────────────────────────────────────────────────────
\usepackage{amsmath}
\numberwithin{equation}{chapter}   % nomor persamaan per-bab: (3.1), (3.2)
\usepackage{amssymb}
\usepackage{siunitx}
\sisetup{output-decimal-marker={,}}

% ── Mikrografi & hyperlink ───────────────────────────────────────────
\usepackage{microtype}
\setlength{\emergencystretch}{3em}
\usepackage[hidelinks]{hyperref}
\hypersetup{
  pdftitle={Judul Tesis Anda},
  pdfauthor={Nama Anda},
}

% ── Warna ────────────────────────────────────────────────────────────
\usepackage{xcolor}
\definecolor{cProses}{HTML}{4A6D8C}    \definecolor{cProsesBg}{HTML}{DCE6F0}
\definecolor{cData}{HTML}{8C7B4A}      \definecolor{cDataBg}{HTML}{F0EBDC}
\definecolor{cKeputusan}{HTML}{8C6A4A} \definecolor{cKeputusanBg}{HTML}{F5E9D8}
\definecolor{cTerm}{HTML}{4A4A4A}      \definecolor{cTermBg}{HTML}{E4E4E4}
\definecolor{cGrup}{HTML}{9AA7B4}      \definecolor{cGrupBg}{HTML}{F2F5F8}
\definecolor{cOk}{HTML}{3E7D5A}        \definecolor{cOkBg}{HTML}{DFF0E6}
\definecolor{cBad}{HTML}{A4503C}       \definecolor{cBadBg}{HTML}{F2DED7}
\definecolor{kodeBg}{HTML}{F7F7F4}
\definecolor{kodeKomentar}{HTML}{5A7D5A}
\definecolor{kodeKunci}{HTML}{4A4A8C}
\definecolor{kodeString}{HTML}{8C5A3C}

% ── TikZ / pgfplots ──────────────────────────────────────────────────
\usepackage{tikz}
\usepackage{pgfplots}
\pgfplotsset{compat=1.18}
\usetikzlibrary{
  shapes.geometric, arrows.meta, positioning,
  calc, fit, backgrounds, decorations.pathreplacing
}

% node styles (flowchart monokrom)
\tikzset{
  term/.style={ellipse, draw=black, line width=.5pt,
    minimum width=26mm, minimum height=9mm,
    align=center, font=\footnotesize},
  proc/.style={rectangle, draw=black, line width=.5pt,
    minimum width=34mm, minimum height=11mm,
    inner xsep=3.4mm, align=center, font=\footnotesize,
    append after command={\pgfextra{
      \draw[line width=.5pt]
        ([xshift=1.6mm]\tikzlastnode.north west)
        -- ([xshift=1.6mm]\tikzlastnode.south west);
      \draw[line width=.5pt]
        ([xshift=-1.6mm]\tikzlastnode.north east)
        -- ([xshift=-1.6mm]\tikzlastnode.south east);}}},
  procplain/.style={rectangle, draw=black, line width=.5pt,
    minimum width=34mm, minimum height=11mm,
    align=center, font=\footnotesize},
  io/.style={trapezium, trapezium left angle=70, trapezium right angle=110,
    draw=black, line width=.5pt,
    minimum width=28mm, minimum height=9mm,
    align=center, inner sep=2pt, font=\footnotesize},
  dec/.style={diamond, aspect=2, draw=black, line width=.5pt,
    align=center, inner sep=1pt, font=\footnotesize},
  arr/.style={-{Stealth[length=2.2mm]}, draw=black, line width=.5pt},
  flbl/.style={font=\footnotesize, fill=white, inner sep=1pt},
}

% ── Listing kode ─────────────────────────────────────────────────────
\usepackage{listings}
\renewcommand{\lstlistingname}{Kode Sumber}
\lstdefinestyle{kode}{
  backgroundcolor=\color{kodeBg},
  basicstyle=\ttfamily\footnotesize,
  commentstyle=\color{kodeKomentar}\itshape,
  keywordstyle=\color{kodeKunci}\bfseries,
  stringstyle=\color{kodeString},
  numberstyle=\ttfamily\scriptsize\color{gray},
  numbers=left, numbersep=8pt, stepnumber=1,
  frame=single, rulecolor=\color{cGrup},
  breaklines=true, breakatwhitespace=false,
  showstringspaces=false, tabsize=2,
  captionpos=t, xleftmargin=18pt, framexleftmargin=14pt,
  literate={→}{$\rightarrow$}{1} {≤}{$\le$}{1} {≥}{$\ge$}{1}
           {×}{$\times$}{1} {—}{{-{}-}}1 {–}{{-}}1
}
\lstset{style=kode}

% ── Makro bantu ──────────────────────────────────────────────────────
\usepackage[export]{adjustbox}
\newcommand{\gambarbesar}[1]{%
  \adjustbox{max width=\textwidth, max totalheight=0.82\textheight, center}{#1}}

% ── Metadata dokumen ─────────────────────────────────────────────────
\title{%
  Judul Tesis\\
  \large Subjudul (jika ada)%
}
\author{Nama Mahasiswa}
\date{\the\year}

% ═════════════════════════════════════════════════════════════════════
\begin{document}

% ── Cover (tidak bernomor halaman) ───────────────────────────────────
\pagenumbering{gobble}
% =====================================================================
%  HALAMAN JUDUL (COVER)
%  Ganti semua teks dalam kurung siku [ ] sesuai data Anda.
% =====================================================================

\begin{titlepage}
\centering

% Logo institusi
\includegraphics[width=3.5cm]{assets/logo_institusi.png}\\[0.5cm]

{\Large\bfseries [NAMA UNIVERSITAS]}\\[0.15cm]
{\large\bfseries [NAMA FAKULTAS]}\\[0.05cm]
{\large\bfseries [NAMA PROGRAM STUDI]}\\[1.5cm]

\rule{\linewidth}{0.5pt}\\[0.3cm]
{\LARGE\bfseries [JUDUL TESIS ANDA]}\\[0.1cm]
{\large [Subjudul jika ada]}\\[0.1cm]
\rule{\linewidth}{0.5pt}\\[1.2cm]

{\normalsize Tesis ini diajukan sebagai salah satu syarat untuk memperoleh gelar}\\[0.1cm]
{\normalsize\bfseries [Sarjana / Magister / Doktor] [Bidang Ilmu]}\\[1.5cm]

{\normalsize Disusun oleh:}\\[0.3cm]
{\large\bfseries [NAMA LENGKAP MAHASISWA]}\\[0.1cm]
{\normalsize NPM: [NOMOR POKOK MAHASISWA]}\\[2cm]

{\normalsize Dosen Pembimbing:}\\[0.2cm]
{\normalsize\bfseries [Nama Pembimbing I, Gelar]}\\[0.05cm]
{\normalsize\bfseries [Nama Pembimbing II, Gelar]}\\[2.5cm]

{\normalsize [KOTA], [BULAN] [TAHUN]}

\end{titlepage}
\clearpage


% ── Front matter (nomor halaman romawi) ──────────────────────────────
\pagenumbering{roman}
\setcounter{page}{1}

% =====================================================================
%  LEMBAR PENGESAHAN
% =====================================================================

\chapter*{Lembar Pengesahan}
\addcontentsline{toc}{chapter}{Lembar Pengesahan}

\begin{center}
{\large\bfseries LEMBAR PENGESAHAN}\\[1cm]

Tesis dengan judul:\\[0.3cm]
{\bfseries ``[JUDUL TESIS ANDA]''}\\[0.3cm]
yang ditulis oleh:\\[0.3cm]
{\bfseries [Nama Mahasiswa]}\\
NPM: [Nomor Pokok Mahasiswa]\\[1cm]

telah diuji dan dinyatakan lulus pada tanggal [Tanggal Ujian].
\end{center}

\vspace{1.5cm}

\noindent
\begin{tabular}{p{6.5cm}p{6.5cm}}
Menyetujui, & \\[0.3cm]
Pembimbing I, & Pembimbing II, \\[2.5cm]
\underline{[Nama Pembimbing I, Gelar]} & \underline{[Nama Pembimbing II, Gelar]} \\
NIP. [XXXXXXXXXX] & NIP. [XXXXXXXXXX] \\
\end{tabular}

\vspace{1.5cm}

\noindent
\begin{tabular}{p{13.5cm}}
Mengetahui, \\[0.3cm]
Ketua Program Studi [Nama Prodi], \\[2.5cm]
\underline{[Nama Kaprodi, Gelar]} \\
NIP. [XXXXXXXXXX] \\
\end{tabular}

\clearpage

% =====================================================================
%  LEMBAR PERNYATAAN KEASLIAN
% =====================================================================

\chapter*{Pernyataan Keaslian}
\addcontentsline{toc}{chapter}{Pernyataan Keaslian}

\noindent Saya yang bertanda tangan di bawah ini:

\vspace{0.5cm}
\begin{tabular}{lll}
Nama            & : & [Nama Lengkap Mahasiswa] \\
NPM             & : & [Nomor Pokok Mahasiswa]  \\
Program Studi   & : & [Nama Program Studi]     \\
Judul Tesis     & : & [Judul Tesis]            \\
\end{tabular}
\vspace{0.5cm}

\noindent menyatakan bahwa tesis ini adalah hasil karya sendiri dan tidak mengandung karya yang
pernah diajukan untuk memperoleh gelar akademik di suatu perguruan tinggi, dan tidak terdapat
karya atau pendapat yang pernah ditulis atau diterbitkan oleh orang lain, kecuali yang secara
tertulis diacu dalam naskah ini dan disebutkan dalam daftar pustaka.

\vspace{2cm}

\noindent
\begin{tabular}{p{8cm}p{5cm}}
 & [Kota], [Tanggal] \\[0.2cm]
 & Yang menyatakan, \\[2.5cm]
 & \underline{[Nama Mahasiswa]} \\
 & NPM. [XXXXXXXXXX] \\
\end{tabular}

\clearpage

% =====================================================================
%  ABSTRAK (Indonesia) + ABSTRACT (English)
% =====================================================================

\chapter*{Abstrak}
\addcontentsline{toc}{chapter}{Abstrak}

\noindent\textbf{[JUDUL TESIS DALAM BAHASA INDONESIA]}

\medskip\noindent\textbf{Nama:} [Nama Mahasiswa] \quad \textbf{NPM:} [XXXXXXXXXX]

\medskip
\noindent
[Isi abstrak dalam Bahasa Indonesia. Panjang 200--300 kata. Abstrak mencakup: (1) latar
belakang dan motivasi penelitian, (2) tujuan penelitian, (3) metode/pendekatan yang
digunakan, (4) hasil utama yang diperoleh, (5) kesimpulan dan kontribusi. Tidak mengandung
rujukan pustaka, singkatan tanpa penjelasan, atau persamaan matematis.]

\medskip
\noindent\textbf{Kata kunci:} [kata kunci 1]; [kata kunci 2]; [kata kunci 3]; [kata kunci 4]; [kata kunci 5]

\clearpage

% ── Abstract (English) ───────────────────────────────────────────────
\chapter*{Abstract}
\addcontentsline{toc}{chapter}{Abstract}

\noindent\textbf{[THESIS TITLE IN ENGLISH]}

\medskip\noindent\textbf{Name:} [Student Name] \quad \textbf{SID:} [XXXXXXXXXX]

\medskip
\noindent
[Abstract content in English. Same structure as the Indonesian abstract: (1) background
and motivation, (2) research objectives, (3) methods/approach, (4) main results,
(5) conclusions and contributions. 200--300 words. No references, unexplained acronyms,
or mathematical equations.]

\medskip
\noindent\textbf{Keywords:} [keyword 1]; [keyword 2]; [keyword 3]; [keyword 4]; [keyword 5]

\clearpage

% =====================================================================
%  KATA PENGANTAR
% =====================================================================

\chapter*{Kata Pengantar}
\addcontentsline{toc}{chapter}{Kata Pengantar}

Puji syukur penulis panjatkan kepada [Tuhan/Allah SWT] atas segala rahmat dan karunia-Nya
sehingga tesis ini dapat diselesaikan.

Tesis berjudul \textit{``[Judul Tesis]''} ini disusun sebagai salah satu syarat untuk
memperoleh gelar [Sarjana/Magister/Doktor] di [Nama Program Studi], [Nama Fakultas],
[Nama Universitas].

Penyelesaian tesis ini tidak terlepas dari bantuan berbagai pihak. Penulis mengucapkan
terima kasih kepada:

\begin{enumerate}
  \item {[Nama Pembimbing I, Gelar]}, selaku dosen pembimbing I yang telah memberikan
        bimbingan, arahan, dan motivasi selama proses penelitian.
  \item {[Nama Pembimbing II, Gelar]}, selaku dosen pembimbing II yang telah memberikan
        masukan berharga.
  \item {[Nama Ketua Prodi, Gelar]}, selaku Ketua Program Studi [Nama Prodi].
  \item Seluruh dosen dan staf [Nama Fakultas] yang telah memberikan ilmu dan dukungan.
  \item Keluarga tercinta atas doa dan dukungan yang tak ternilai.
  \item Rekan-rekan [angkatan/lab/kelompok] atas kebersamaan dan semangat yang diberikan.
\end{enumerate}

Penulis menyadari bahwa tesis ini masih jauh dari sempurna. Oleh karena itu, kritik dan
saran yang membangun sangat diharapkan demi penyempurnaan tulisan ini.

Semoga tesis ini bermanfaat bagi pembaca dan perkembangan ilmu pengetahuan.

\vspace{1cm}

\noindent
\begin{tabular}{p{8cm}p{5cm}}
 & [Kota], [Bulan] [Tahun] \\[0.2cm]
 & Penulis, \\[2.5cm]
 & [Nama Mahasiswa] \\
\end{tabular}

\clearpage


\tableofcontents
\listoffigures
\listoftables

\clearpage
\pagenumbering{arabic}
\setcounter{page}{1}

% ── Bab-bab utama ────────────────────────────────────────────────────
% =====================================================================
%  BAB 1 — PENDAHULUAN
% =====================================================================

\chapter{Pendahuluan}

\section{Latar Belakang}
\label{sec:1.1}

[Uraikan latar belakang masalah secara komprehensif. Mulai dari kondisi umum, identifikasi
kesenjangan (\emph{gap}) yang ada, hingga urgensi penelitian ini dilakukan. Dukung dengan
data dan referensi empiris. Panjang: 3--5 paragraf.]

[Paragraf pertama: konteks besar / kondisi terkini di bidang ini.]

[Paragraf kedua: masalah spesifik yang belum terpecahkan / keterbatasan pendekatan yang ada.]

[Paragraf ketiga: solusi yang diusulkan secara ringkas dan mengapa pendekatan ini relevan.]

\section{Rumusan Masalah}
\label{sec:1.2}

Berdasarkan latar belakang di atas, penelitian ini merumuskan permasalahan sebagai berikut:

\begin{enumerate}
  \item Bagaimana [rumusan masalah pertama]?
  \item Bagaimana [rumusan masalah kedua]?
  \item Seberapa [rumusan masalah ketiga, biasanya terkait kinerja/evaluasi]?
\end{enumerate}

\section{Tujuan Penelitian}
\label{sec:1.3}

Penelitian ini bertujuan untuk:

\begin{enumerate}
  \item [Tujuan pertama, sesuai dengan rumusan masalah 1.]
  \item [Tujuan kedua, sesuai dengan rumusan masalah 2.]
  \item [Tujuan ketiga, sesuai dengan rumusan masalah 3.]
\end{enumerate}

\section{Batasan Masalah}
\label{sec:1.4}

Agar penelitian terfokus, ditetapkan batasan-batasan sebagai berikut:

\begin{enumerate}
  \item [Batasan 1: lingkup teknologi, dataset, atau platform yang digunakan.]
  \item [Batasan 2: jenis eksperimen atau skenario yang dicakup.]
  \item [Batasan 3: asumsi yang disederhanakan dari kondisi nyata.]
\end{enumerate}

\section{Manfaat Penelitian}
\label{sec:1.5}

\textbf{Manfaat Teoritis:}
[Jelaskan kontribusi terhadap pengembangan ilmu pengetahuan, misalnya: memperkaya literatur
di bidang X, menyediakan kerangka referensi untuk penelitian sejenis, dsb.]

\medskip
\textbf{Manfaat Praktis:}
[Jelaskan manfaat langsung bagi praktisi atau industri, misalnya: membantu tim keamanan
dalam Y, menyederhanakan proses Z, dsb.]

\section{Sistematika Penulisan}
\label{sec:1.6}

Tesis ini disusun dalam tujuh bab dengan sistematika sebagai berikut.

\textbf{Bab~1 Pendahuluan} menyajikan latar belakang, rumusan masalah, tujuan, batasan,
dan manfaat penelitian.

\textbf{Bab~2 Tinjauan Pustaka} membahas teori-teori dasar dan penelitian terdahulu yang
relevan sebagai landasan penelitian ini.

\textbf{Bab~3 Metodologi Penelitian} menguraikan desain penelitian, metode pengumpulan data,
dan teknik analisis yang digunakan.

\textbf{Bab~4 Perancangan Sistem} mendeskripsikan arsitektur dan desain komponen sistem yang
dibangun.

\textbf{Bab~5 Implementasi} menjelaskan realisasi sistem berdasarkan perancangan pada Bab~4.

\textbf{Bab~6 Pengujian dan Analisis} menyajikan hasil eksperimen beserta analisis dan
pembahasan secara mendalam.

\textbf{Bab~7 Kesimpulan dan Saran} merangkum temuan utama penelitian dan mengusulkan arah
penelitian lanjutan.

% =====================================================================
%  BAB 2 — TINJAUAN PUSTAKA
% =====================================================================

\chapter{Tinjauan Pustaka}

% ── 2.1 Landasan Teori ───────────────────────────────────────────────
\section{Landasan Teori}
\label{sec:2.1}

\subsection{[Konsep Utama 1]}
\label{sec:2.1.1}

[Penjelasan konsep utama pertama yang mendasari penelitian ini. Definisi, cara kerja,
dan relevansinya terhadap penelitian. Sertakan rujukan ke literatur primer.]

\subsection{[Konsep Utama 2]}
\label{sec:2.1.2}

[Penjelasan konsep utama kedua.]

\subsection{[Konsep Utama 3]}
\label{sec:2.1.3}

[Penjelasan konsep utama ketiga.]

% Contoh persamaan matematis
% Misalnya: F1-score
\begin{equation}
  F_1 = 2 \cdot \frac{\text{Presisi} \times \text{Recall}}{\text{Presisi} + \text{Recall}}
  \label{eq:f1}
\end{equation}

% Contoh tabel perbandingan
% \begin{table}[H]
%   \centering
%   \caption{Perbandingan [hal yang dibandingkan]}
%   \label{tab:2.perbandingan}
%   \begin{tabular}{llll}
%     \toprule
%     \textbf{Metode} & \textbf{Kelebihan} & \textbf{Kekurangan} & \textbf{Sumber} \\
%     \midrule
%     {[Metode A]} & [+] & [-] & [Penulis, Tahun] \\
%     {[Metode B]} & [+] & [-] & [Penulis, Tahun] \\
%     \bottomrule
%   \end{tabular}
% \end{table}

% ── 2.2 Penelitian Terdahulu ─────────────────────────────────────────
\section{Penelitian Terdahulu}
\label{sec:2.2}

[Ulas penelitian-penelitian yang relevan secara kritis. Identifikasi kontribusi, metode,
dataset, dan hasil dari setiap penelitian. Gunakan paragraf tematis, bukan hanya daftar
rangkuman.

Struktur yang baik: kelompokkan penelitian berdasarkan pendekatan, bukan kronologi.
Tunjukkan bagaimana setiap penelitian berkaitan dengan gap yang Anda isi.]

\subsection{[Kelompok Penelitian 1: Pendekatan X]}
\label{sec:2.2.1}

[Ulas beberapa penelitian yang menggunakan pendekatan X. Sintesis persamaan dan
perbedaannya.]

\subsection{[Kelompok Penelitian 2: Pendekatan Y]}
\label{sec:2.2.2}

[Ulas pendekatan Y dan kelemahannya yang menjadi motivasi penelitian Anda.]

% ── 2.3 Posisi Penelitian ────────────────────────────────────────────
\section{Posisi Penelitian}
\label{sec:2.3}

Berdasarkan tinjauan terhadap penelitian-penelitian yang telah diuraikan, Tabel~\ref{tab:2.posisi}
merangkum posisi penelitian ini dibandingkan karya terdahulu.

\begin{table}[H]
  \centering
  \caption{Posisi penelitian ini terhadap penelitian terdahulu}
  \label{tab:2.posisi}
  \begin{tabular}{p{2.9cm} p{1.8cm} p{1.8cm} p{1.8cm} p{1.8cm} p{1.8cm}}
    \toprule
    \textbf{Penelitian} & \textbf{[Aspek 1]} & \textbf{[Aspek 2]} & \textbf{[Aspek 3]}
                        & \textbf{[Aspek 4]} & \textbf{[Aspek 5]} \\
    \midrule
    {[Peneliti A, Tahun]} & \checkmark & -- & \checkmark & -- & -- \\
    {[Peneliti B, Tahun]} & \checkmark & \checkmark & -- & -- & -- \\
    \textbf{Penelitian ini} & \checkmark & \checkmark & \checkmark & \checkmark & \checkmark \\
    \bottomrule
  \end{tabular}
\end{table}

[Paragraf yang menjelaskan kebaruan (\emph{novelty}) penelitian ini berdasarkan tabel di atas.
Apa yang belum dilakukan orang lain dan yang Anda lakukan?]

% =====================================================================
%  BAB 3 — METODOLOGI PENELITIAN
% =====================================================================

\chapter{Metodologi Penelitian}

\section{Desain Penelitian}
\label{sec:3.1}

Penelitian ini menggunakan pendekatan [eksperimental / deskriptif / kuantitatif / kualitatif / campuran].
[Jelaskan alasan pemilihan desain penelitian ini dan bagaimana desain tersebut menjawab
rumusan masalah yang telah ditetapkan.]

% Contoh gambar alur penelitian
\begin{figure}[H]
  \centering
  \begin{tikzpicture}[node distance=1cm]
    \node[term]    (mulai)   {Mulai};
    \node[proc]    (studi)   [below=of mulai]   {Studi Literatur};
    \node[proc]    (rumus)   [below=of studi]   {Perumusan Masalah};
    \node[proc]    (data)    [below=of rumus]   {Pengumpulan Data};
    \node[proc]    (analisis)[below=of data]    {Analisis \& Perancangan};
    \node[proc]    (impl)    [below=of analisis]{Implementasi};
    \node[proc]    (uji)     [below=of impl]    {Pengujian \& Evaluasi};
    \node[dec]     (valid)   [below=of uji]     {Tujuan \\ tercapai?};
    \node[proc]    (perbaiki)[right=2cm of valid]{Perbaikan};
    \node[proc]    (laporan) [below=of valid]   {Penulisan Laporan};
    \node[term]    (selesai) [below=of laporan] {Selesai};

    \draw[arr] (mulai)    -- (studi);
    \draw[arr] (studi)    -- (rumus);
    \draw[arr] (rumus)    -- (data);
    \draw[arr] (data)     -- (analisis);
    \draw[arr] (analisis) -- (impl);
    \draw[arr] (impl)     -- (uji);
    \draw[arr] (uji)      -- (valid);
    \draw[arr] (valid)    -- node[flbl,right]{Ya}  (laporan);
    \draw[arr] (valid)    -- node[flbl,above]{Tidak}(perbaiki);
    \draw[arr] (perbaiki) |- (impl);
    \draw[arr] (laporan)  -- (selesai);
  \end{tikzpicture}
  \caption{Alur penelitian}
  \label{fig:3.alur}
\end{figure}

\section{Dataset dan Lingkungan Eksperimen}
\label{sec:3.2}

\subsection{Dataset}
\label{sec:3.2.1}

[Deskripsikan dataset yang digunakan: sumber, ukuran, karakteristik distribusi kelas,
format, dan cara mendapatkannya. Jelaskan apakah dataset publik atau dikumpulkan sendiri.]

\begin{table}[H]
  \centering
  \caption{Karakteristik dataset yang digunakan}
  \label{tab:3.dataset}
  \begin{tabular}{lrrl}
    \toprule
    \textbf{Subset} & \textbf{Jumlah sampel} & \textbf{Proporsi (\%)} & \textbf{Keterangan} \\
    \midrule
    {[Kelas A / Benign]} & [X.XXX] & [XX,X] & [Deskripsi singkat] \\
    {[Kelas B / Attack]} & [X.XXX] & [XX,X] & [Deskripsi singkat] \\
    \midrule
    \textbf{Total}     & [X.XXX] & 100,0  & \\
    \bottomrule
  \end{tabular}
\end{table}

\subsection{Lingkungan Eksperimen}
\label{sec:3.2.2}

Eksperimen dilakukan pada lingkungan sebagai berikut:

\begin{table}[H]
  \centering
  \caption{Spesifikasi lingkungan eksperimen}
  \label{tab:3.lingkungan}
  \begin{tabular}{ll}
    \toprule
    \textbf{Komponen} & \textbf{Spesifikasi} \\
    \midrule
    Sistem Operasi   & [Nama OS, Versi] \\
    Prosesor         & [Nama CPU, Jumlah Core, Frekuensi] \\
    Memori (RAM)     & [Kapasitas GB] \\
    Penyimpanan      & [Tipe, Kapasitas] \\
    Bahasa Pemrograman & [Bahasa, Versi] \\
    Framework/Library& [Nama, Versi] \\
    \bottomrule
  \end{tabular}
\end{table}

\section{Metode yang Digunakan}
\label{sec:3.3}

\subsection{[Metode 1]}
\label{sec:3.3.1}

[Uraikan metode/algoritma pertama yang digunakan. Jelaskan cara kerja, parameter kunci,
dan alasan pemilihan metode ini. Sertakan pseudocode jika perlu.]

% Contoh pseudocode
% \begin{lstlisting}[language=Python, caption={Pseudocode [Metode 1]}, label={lst:3.pseudo}]
% def metode(input):
%     # Langkah 1: inisialisasi
%     hasil = []
%     # Langkah 2: proses
%     for item in input:
%         hasil.append(proses(item))
%     return hasil
% \end{lstlisting}

\subsection{[Metode 2]}
\label{sec:3.3.2}

[Uraikan metode kedua.]

\section{Metrik Evaluasi}
\label{sec:3.4}

Kinerja sistem dievaluasi menggunakan metrik-metrik berikut.

\textbf{Presisi} (\emph{Precision}) mengukur proporsi prediksi positif yang benar:
\begin{equation}
  \text{Presisi} = \frac{TP}{TP + FP}
  \label{eq:presisi}
\end{equation}

\textbf{Recall} (\emph{Sensitivitas}) mengukur kemampuan model mendeteksi kasus positif:
\begin{equation}
  \text{Recall} = \frac{TP}{TP + FN}
  \label{eq:recall}
\end{equation}

\textbf{F1-score} merupakan rata-rata harmonis antara presisi dan recall:
\begin{equation}
  F_1 = 2 \cdot \frac{\text{Presisi} \times \text{Recall}}{\text{Presisi} + \text{Recall}}
  \label{eq:f1}
\end{equation}

\noindent
di mana $TP$ = \emph{True Positive}, $FP$ = \emph{False Positive}, $FN$ = \emph{False Negative}.

[Tambahkan metrik lain yang relevan, misalnya latensi, throughput, AUC-ROC, dsb.]

\section{Prosedur Eksperimen}
\label{sec:3.5}

[Jelaskan langkah-langkah eksperimen secara runtut: pembagian data (train/val/test),
metode validasi (k-fold cross-validation, hold-out), skenario ablasi, dsb. Pastikan
eksperimen dapat direproduksi oleh peneliti lain.]

Eksperimen dilaksanakan dalam tahapan berikut:

\begin{enumerate}
  \item \textbf{Praproses data}: [normalisasi, penanganan nilai hilang, encoding, dsb.]
  \item \textbf{Pelatihan model}: [pembagian data, hiperparameter, strategi tuning.]
  \item \textbf{Pengujian}: [skenario pengujian, baseline yang dibandingkan.]
  \item \textbf{Analisis}: [metrik yang dihitung, visualisasi yang dibuat.]
\end{enumerate}

% =====================================================================
%  BAB 4 — PERANCANGAN SISTEM
% =====================================================================

\chapter{Perancangan Sistem}

\section{Gambaran Umum Sistem}
\label{sec:4.1}

[Berikan gambaran arsitektur sistem secara keseluruhan. Jelaskan komponen-komponen utama,
bagaimana komponen-komponen tersebut berinteraksi, dan bagaimana sistem secara keseluruhan
menjawab rumusan masalah yang ditetapkan.]

\begin{figure}[H]
  \centering
  \gambarbesar{%
    \begin{tikzpicture}[node distance=1.5cm and 2cm]
      % ── Ganti dengan diagram arsitektur Anda ─────────────────────
      \node[procplain, fill=cProsesBg, draw=cProses, text width=3cm]
        (komp1) {Komponen 1};
      \node[procplain, fill=cDataBg, draw=cData, text width=3cm,
            right=of komp1]
        (komp2) {Komponen 2};
      \node[procplain, fill=cOkBg, draw=cOk, text width=3cm,
            right=of komp2]
        (komp3) {Komponen 3};

      \draw[arr] (komp1) -- node[above,font=\scriptsize]{Data} (komp2);
      \draw[arr] (komp2) -- node[above,font=\scriptsize]{Hasil} (komp3);
    \end{tikzpicture}
  }
  \caption{Arsitektur sistem secara keseluruhan}
  \label{fig:4.arsitektur}
\end{figure}

[Deskripsi singkat setiap komponen: peran, input, output.]

\section{Perancangan [Komponen 1]}
\label{sec:4.2}

[Deskripsikan desain komponen pertama secara rinci. Jelaskan algoritma, struktur data,
antarmuka (interface), dan keputusan desain yang diambil beserta alasannya.]

\subsection{[Sub-komponen 1.1]}
\label{sec:4.2.1}

[Detail sub-komponen.]

\subsection{[Sub-komponen 1.2]}
\label{sec:4.2.2}

[Detail sub-komponen.]

\section{Perancangan [Komponen 2]}
\label{sec:4.3}

[Deskripsikan desain komponen kedua.]

\section{Perancangan [Komponen 3]}
\label{sec:4.4}

[Deskripsikan desain komponen ketiga.]

\section{Rancangan Antarmuka Data}
\label{sec:4.5}

[Jelaskan format data yang mengalir antar komponen, skema database (jika ada), format
file input/output, dan protokol komunikasi antar-komponen.]

% Contoh format output JSON
% \begin{lstlisting}[language=,caption={Format data antar-komponen},label={lst:4.format}]
% {
%   "id": "string",
%   "timestamp": "ISO8601",
%   "fitur_a": float,
%   "fitur_b": float,
%   "label": "string"
% }
% \end{lstlisting}

\section{Rancangan Pengujian}
\label{sec:4.6}

[Deskripsikan skenario pengujian yang akan dilakukan pada Bab~6. Apa yang akan diuji,
bagaimana cara mengukurnya, dan kriteria keberhasilan yang ditetapkan.]

\begin{table}[H]
  \centering
  \caption{Skenario pengujian}
  \label{tab:4.skenario}
  \begin{tabular}{clll}
    \toprule
    \textbf{No.} & \textbf{Skenario} & \textbf{Variabel} & \textbf{Metrik} \\
    \midrule
    S1 & [Nama skenario 1] & [Variabel yang diubah] & [Metrik yang diukur] \\
    S2 & [Nama skenario 2] & [Variabel yang diubah] & [Metrik yang diukur] \\
    S3 & [Nama skenario 3] & [Variabel yang diubah] & [Metrik yang diukur] \\
    \bottomrule
  \end{tabular}
\end{table}

% =====================================================================
%  BAB 5 — IMPLEMENTASI
% =====================================================================

\chapter{Implementasi}

\section{Lingkungan Implementasi}
\label{sec:5.1}

Implementasi dilakukan pada lingkungan yang telah dideskripsikan pada
Tabel~\ref{tab:3.lingkungan} di Bab~3. [Tambahkan detail khusus implementasi jika ada
perbedaan atau tambahan dari yang sudah disebutkan.]

\section{Implementasi [Komponen 1]}
\label{sec:5.2}

[Jelaskan bagaimana Komponen 1 diimplementasikan. Hubungkan implementasi ke rancangan
di Bab~4. Sorot keputusan teknis yang diambil selama implementasi, terutama yang berbeda
dari rancangan awal beserta alasannya.]

\subsection{[Modul / Fungsi Kunci 1]}
\label{sec:5.2.1}

[Penjelasan modul. Sertakan snippet kode untuk bagian yang non-trivial atau inovatif.
Jangan tampilkan kode yang sudah jelas; fokus pada bagian yang bernilai penjelasan.]

\begin{lstlisting}[language=Python, caption={[Deskripsi Kode Sumber 1]}, label={lst:5.kode1}]
# Contoh implementasi [fungsi kunci]
def fungsi_kunci(parameter_a, parameter_b):
    """[Docstring singkat.]"""
    # [Langkah 1: jelaskan logika non-trivial di sini]
    hasil_antara = proses_awal(parameter_a)

    # [Langkah 2]
    for item in hasil_antara:
        item = transformasi(item, parameter_b)

    return hasil_antara
\end{lstlisting}

\subsection{[Modul / Fungsi Kunci 2]}
\label{sec:5.2.2}

[Penjelasan modul berikutnya.]

\section{Implementasi [Komponen 2]}
\label{sec:5.3}

[Implementasi komponen kedua.]

\section{Implementasi [Komponen 3]}
\label{sec:5.4}

[Implementasi komponen ketiga.]

\section{Integrasi Antar-Komponen}
\label{sec:5.5}

[Jelaskan bagaimana komponen-komponen dihubungkan satu sama lain. Pipeline data, antrian
pesan, API, konfigurasi, dsb. Tampilkan diagram aliran data atau topologi deployment
jika membantu.]

\begin{figure}[H]
  \centering
  \begin{tikzpicture}[node distance=1.2cm and 2.5cm]
    % ── Ganti dengan diagram integrasi Anda ──────────────────────────
    \node[procplain, fill=cProsesBg, draw=cProses] (a) {[Komponen A]};
    \node[broker,    fill=cDataBg,   draw=cData,
          right=of a] (b) {[Broker/\\Queue]};
    \node[procplain, fill=cOkBg,    draw=cOk,
          right=of b] (c) {[Komponen C]};

    \draw[arr] (a) -- node[above,font=\scriptsize]{[data]} (b);
    \draw[arr] (b) -- node[above,font=\scriptsize]{[event]} (c);
  \end{tikzpicture}
  \caption{Topologi integrasi antar-komponen}
  \label{fig:5.integrasi}
\end{figure}

\section{Kendala Implementasi}
\label{sec:5.6}

[Uraikan kendala teknis yang ditemui selama implementasi dan bagaimana cara mengatasinya.
Ini adalah bagian yang sering dilewatkan tapi bernilai tinggi bagi pembaca.]

\begin{table}[H]
  \centering
  \caption{Kendala implementasi dan solusinya}
  \label{tab:5.kendala}
  \begin{tabular}{p{5cm}p{4cm}p{4cm}}
    \toprule
    \textbf{Kendala} & \textbf{Dampak} & \textbf{Solusi} \\
    \midrule
    [Deskripsi kendala 1] & [Dampak pada sistem] & [Solusi yang diterapkan] \\
    [Deskripsi kendala 2] & [Dampak pada sistem] & [Solusi yang diterapkan] \\
    \bottomrule
  \end{tabular}
\end{table}

% =====================================================================
%  BAB 6 — PENGUJIAN DAN ANALISIS
% =====================================================================

\chapter{Pengujian dan Analisis}

\section{Prosedur Pengujian}
\label{sec:6.1}

[Jelaskan prosedur pengujian yang dilaksanakan, mengacu pada rancangan pengujian di
Seksyen~\ref{sec:4.6}. Uraikan urutan langkah, konfigurasi yang digunakan, dan cara
hasil dikumpulkan.]

% Contoh: diagram alur pengujian
\begin{figure}[H]
  \centering
  \begin{tikzpicture}[node distance=0.9cm]
    \node[term]     (mulai)  {Mulai};
    \node[io]       (setup)  [below=of mulai]  {Inisialisasi lingkungan};
    \node[proc]     (data)   [below=of setup]  {Muat data uji};
    \node[proc]     (jalankan)[below=of data]  {Jalankan skenario};
    \node[io]       (kumpul) [below=of jalankan]{Kumpulkan metrik};
    \node[dec]      (ulang)  [below=of kumpul] {Semua skenario\\selesai?};
    \node[proc]     (analisis)[below=of ulang] {Analisis \& visualisasi};
    \node[term]     (selesai)[below=of analisis]{Selesai};
    \node[proc]     (iterasi)[right=2cm of ulang]{Skenario\\berikutnya};

    \draw[arr] (mulai)    -- (setup);
    \draw[arr] (setup)    -- (data);
    \draw[arr] (data)     -- (jalankan);
    \draw[arr] (jalankan) -- (kumpul);
    \draw[arr] (kumpul)   -- (ulang);
    \draw[arr] (ulang)    -- node[flbl,right]{Ya} (analisis);
    \draw[arr] (ulang)    -- node[flbl,above]{Tidak} (iterasi);
    \draw[arr] (iterasi)  |- (data);
    \draw[arr] (analisis) -- (selesai);
  \end{tikzpicture}
  \caption{Alur prosedur pengujian}
  \label{fig:6.alur}
\end{figure}

% ── 6.2 Hasil Pengujian Utama ────────────────────────────────────────
\section{Hasil Pengujian [Aspek Utama 1]}
\label{sec:6.2}

[Sajikan hasil pengujian skenario pertama. Gunakan tabel untuk angka, grafik untuk tren.
Deskripsikan temuan secara faktual sebelum memberikan interpretasi.]

\begin{table}[H]
  \centering
  \caption{Hasil [Aspek Utama 1] pada berbagai konfigurasi}
  \label{tab:6.hasil1}
  \begin{tabular}{lrrrr}
    \toprule
    \textbf{Konfigurasi} & \textbf{Presisi} & \textbf{Recall} & \textbf{F1} & \textbf{Latensi (ms)} \\
    \midrule
    {[Konfigurasi A]} & [0,XX] & [0,XX] & [0,XX] & [X,XX] \\
    {[Konfigurasi B]} & [0,XX] & [0,XX] & [0,XX] & [X,XX] \\
    {[Konfigurasi C]} & [0,XX] & [0,XX] & [0,XX] & [X,XX] \\
    \textbf{Sistem ini} & \textbf{[0,XX]} & \textbf{[0,XX]} & \textbf{[0,XX]} & \textbf{[X,XX]} \\
    \bottomrule
  \end{tabular}
\end{table}

% Contoh grafik pgfplots (bar chart)
\begin{figure}[H]
  \centering
  \begin{tikzpicture}
    \begin{axis}[
      ybar, bar width=0.5cm,
      width=0.85\textwidth, height=6cm,
      ylabel={F1-score},
      xtick=data,
      xticklabels={[Konf A],[Konf B],[Konf C],[Sistem ini]},
      ymin=0, ymax=1.05,
      ymajorgrids=true, grid style=dashed,
      legend pos=north west,
      tick label style={font=\small},
      label style={font=\small},
    ]
      \addplot coordinates {(0,0.80) (1,0.85) (2,0.88) (3,0.95)};
    \end{axis}
  \end{tikzpicture}
  \caption{Perbandingan F1-score antar konfigurasi}
  \label{fig:6.f1}
\end{figure}

[Analisis: mengapa konfigurasi terbaik mengungguli yang lain? Faktor apa yang berkontribusi?]

% ── 6.3 Hasil Pengujian Aspek 2 ──────────────────────────────────────
\section{Hasil Pengujian [Aspek 2]}
\label{sec:6.3}

[Sajikan hasil skenario kedua. Contoh: pengujian skalabilitas, latensi, throughput, dsb.]

% Contoh: grafik garis (latensi vs beban)
\begin{figure}[H]
  \centering
  \begin{tikzpicture}
    \begin{axis}[
      width=0.85\textwidth, height=6cm,
      xlabel={Beban (\emph{request/s})},
      ylabel={Latensi median (ms)},
      xmin=0, ymin=0,
      grid=both, grid style={line width=.2pt, draw=gray!30},
      major grid style={line width=.4pt, draw=gray!60},
      legend pos=north west,
      tick label style={font=\small},
      label style={font=\small},
    ]
      \addplot[mark=o, thick] coordinates {
        (100,1.2) (500,1.5) (1000,2.1) (2000,3.8) (5000,9.5)
      };
      \addlegendentry{[Sistem ini]}
      \addplot[mark=square, thick, dashed] coordinates {
        (100,2.0) (500,3.2) (1000,5.8) (2000,12.0) (5000,35.0)
      };
      \addlegendentry{[Baseline]}
    \end{axis}
  \end{tikzpicture}
  \caption{Latensi median terhadap beban}
  \label{fig:6.latensi}
\end{figure}

% ── 6.4 Analisis Ablasi ──────────────────────────────────────────────
\section{Analisis Ablasi}
\label{sec:6.4}

[Tunjukkan kontribusi masing-masing komponen/fitur dengan mematikannya satu per satu
dan mengukur dampak performa. Ini adalah bagian yang membuktikan bahwa setiap pilihan
desain berkontribusi nyata.]

\begin{table}[H]
  \centering
  \caption{Hasil ablasi komponen sistem}
  \label{tab:6.ablasi}
  \begin{tabular}{lcccr}
    \toprule
    \textbf{Variasi} & \textbf{[Komponen A]} & \textbf{[Komponen B]} & \textbf{[Komponen C]}
                     & \textbf{F1} \\
    \midrule
    Sistem penuh     & \checkmark & \checkmark & \checkmark & [0,XX] \\
    Tanpa [Komp C]   & \checkmark & \checkmark & --         & [0,XX] \\
    Tanpa [Komp B]   & \checkmark & --         & \checkmark & [0,XX] \\
    Tanpa [Komp A]   & --         & \checkmark & \checkmark & [0,XX] \\
    \emph{Baseline}  & --         & --         & --         & [0,XX] \\
    \bottomrule
  \end{tabular}
\end{table}

[Diskusi: komponen mana yang paling berkontribusi dan mengapa.]

% ── 6.5 Perbandingan dengan Penelitian Terdahulu ─────────────────────
\section{Perbandingan dengan Penelitian Terdahulu}
\label{sec:6.5}

[Bandingkan hasil sistem Anda dengan hasil yang dilaporkan pada penelitian terdahulu
(Bab~2). Perhatikan: perbandingan yang adil mensyaratkan dataset dan skenario yang sama
atau setara. Jelaskan keterbatasan perbandingan jika ada.]

\begin{table}[H]
  \centering
  \caption{Perbandingan dengan penelitian terdahulu}
  \label{tab:6.komparasi}
  \begin{tabular}{lp{2.8cm}rrr}
    \toprule
    \textbf{Penelitian} & \textbf{Dataset} & \textbf{F1} & \textbf{Latensi (ms)}
                        & \textbf{[Metrik lain]} \\
    \midrule
    {[Peneliti A, Tahun]} & [Nama dataset] & [0,XX] & [X,X] & [X] \\
    {[Peneliti B, Tahun]} & [Nama dataset] & [0,XX] & [X,X] & [X] \\
    \textbf{Penelitian ini} & [Nama dataset] & \textbf{[0,XX]} & \textbf{[X,X]} & \textbf{[X]} \\
    \bottomrule
  \end{tabular}
\end{table}

% ── 6.6 Diskusi ──────────────────────────────────────────────────────
\section{Diskusi}
\label{sec:6.6}

\subsection{Keterbatasan Sistem}
\label{sec:6.6.1}

[Bahas secara jujur keterbatasan sistem yang dibangun: kondisi di mana performa menurun,
asumsi yang mungkin tidak berlaku di lingkungan produksi nyata, dsb.]

\subsection{Ancaman terhadap Validitas}
\label{sec:6.6.2}

[Identifikasi ancaman terhadap validitas hasil eksperimen: \emph{internal validity}
(apakah perubahan memang disebabkan sistem Anda?), \emph{external validity} (apakah
hasil dapat digeneralisasi?), dan \emph{construct validity} (apakah metrik yang dipilih
benar-benar mengukur yang dimaksud?).]

% =====================================================================
%  BAB 7 — KESIMPULAN DAN SARAN
% =====================================================================

\chapter{Kesimpulan dan Saran}

\section{Kesimpulan}
\label{sec:7.1}

Penelitian ini telah berhasil [menjawab/merancang/membangun/membuktikan] [tujuan utama
penelitian]. Berdasarkan hasil pengujian dan analisis yang telah dilakukan, dapat
ditarik kesimpulan sebagai berikut.

\begin{enumerate}
  \item [Kesimpulan 1 — jawaban langsung atas Rumusan Masalah 1. Sertakan angka/fakta
        konkret dari hasil pengujian, bukan pernyataan umum.]

  \item [Kesimpulan 2 — jawaban atas Rumusan Masalah 2. Spesifik dan terukur.]

  \item [Kesimpulan 3 — jawaban atas Rumusan Masalah 3. Bandingkan dengan baseline atau
        penelitian terdahulu jika relevan.]
\end{enumerate}

[Paragraf penutup yang merangkum kontribusi utama penelitian dan signifikansinya terhadap
bidang keilmuan yang bersangkutan.]

\section{Saran}
\label{sec:7.2}

Berdasarkan temuan dan keterbatasan yang diidentifikasi, beberapa saran untuk penelitian
dan pengembangan lebih lanjut adalah sebagai berikut.

\begin{enumerate}
  \item \textbf{[Saran 1: Perbaikan teknis yang spesifik.]}
        [Jelaskan dengan singkat mengapa ini penting dan bagaimana bisa dilakukan.]

  \item \textbf{[Saran 2: Perluasan cakupan eksperimen.]}
        [Misalnya: menguji pada dataset yang lebih beragam, lingkungan produksi nyata,
        atau skenario serangan yang lebih kompleks.]

  \item \textbf{[Saran 3: Arah penelitian lanjutan.]}
        [Pertanyaan terbuka yang muncul dari penelitian ini yang layak dieksplorasi
        lebih jauh.]

  \item \textbf{[Saran 4: Aspek yang belum dicakup dalam penelitian ini.]}
        [Misalnya: integrasi dengan sistem yang lebih besar, evaluasi biaya operasional,
        atau aspek privasi dan etika yang relevan.]
\end{enumerate}


% ── Daftar pustaka ───────────────────────────────────────────────────
% =====================================================================
%  DAFTAR PUSTAKA
%  Gaya: Harvard Anglia Ruskin University (ARU)
%  Format manual — tanpa BibTeX
%  Urutan: alfabet nama belakang penulis pertama
% =====================================================================

\chapter*{Daftar Pustaka}
\addcontentsline{toc}{chapter}{Daftar Pustaka}

\begingroup
\setlength{\parindent}{-0.5cm}
\setlength{\leftskip}{0.5cm}
\setlength{\parskip}{6pt}

% ── FORMAT HARVARD ARU ────────────────────────────────────────────────
%
% JURNAL
%   Nama Belakang, Inisial. (Tahun) 'Judul artikel',
%   \textit{Nama Jurnal}, Volume(Nomor), hlm. XX--XX.
%   doi: xxxxx [Diakses: DD Bulan YYYY].
%
% BUKU
%   Nama Belakang, Inisial. (Tahun) \textit{Judul Buku}.
%   Edisi ke-X. Kota: Penerbit.
%
% BAB DALAM BUKU YANG DIEDIT
%   Nama Belakang, Inisial. (Tahun) 'Judul bab', dalam
%   Nama Belakang Editor, Inisial. (ed./eds.)
%   \textit{Judul Buku}. Kota: Penerbit, hlm. XX--XX.
%
% PROSIDING KONFERENSI
%   Nama Belakang, Inisial. (Tahun) 'Judul makalah',
%   \textit{Nama Konferensi}, Kota, DD--DD Bulan YYYY.
%   Kota: Penerbit, hlm. XX--XX.
%
% SITUS WEB / DOKUMEN ONLINE
%   Nama Belakang, Inisial. / Nama Organisasi (Tahun)
%   \textit{Judul Halaman}. Tersedia di: URL
%   [Diakses: DD Bulan YYYY].
%
% DUA ATAU TIGA PENULIS: pisah dengan 'dan'
%   Doe, J. dan Smith, A. (Tahun) ...
%
% EMPAT PENULIS ATAU LEBIH: tulis semua sampai penulis terakhir
%   (tidak disingkat \textit{et al.} dalam daftar pustaka —
%    \textit{et al.} hanya untuk sitasi dalam teks)
% ─────────────────────────────────────────────────────────────────────

% Contoh jurnal:
[Nama Belakang], [I]. ([Tahun]) '[Judul artikel]',
\textit{[Nama Jurnal]}, [Vol]([No]), hlm.~[XX--XX].
doi:~[xxxxx] [Diakses: DD Bulan YYYY].

% Contoh buku:
[Nama Belakang], [I]. ([Tahun]) \textit{[Judul Buku]}.
[Kota]: [Penerbit].

% Tambahkan referensi di bawah, diurutkan alfabet nama belakang penulis pertama.

\endgroup


\end{document}
